% Данный файл распространяется под лицензией CC BY 4.0. Текст лицензии размещён на https://creativecommons.org/licenses/by/4.0/
% (c) Кафедра системного программирования, 2026

% Компилировать XeTeX

% Содержательная часть отзыва: поля, которые надо заполнить

% Тип практики, вставляется во фразу "отчёт ... практике". Бывает учебная, производственная, преддипломная, педагогическая.
\def\practiceType{об учебной}

% ФИО (полное) студента
\def\studentName{Фамилия Имя Отчество В Родительном Падеже}

% тема практики как на титульнике
\def\practiceTitle{Тема учебной практики}

% Семестр, за который практика
\def\practiceSemester{второй}

% Оценки по критериям

% Оценка по критерию "Цель практики достигнута полностью", от 1 (совсем нет) до 5 (абсолютно да)
\def\completenessGrade{5}

% Комментарий по критерию в свободной форме (оставьте пустым, если нет)
\def\completenessComment{}

% Оценка по критерию "Этапы работы выполнялись в ожидаемые сроки, работа велась непрерывно в течение семестра", от 1 (совсем нет) до 5 (абсолютно да)
\def\continuityGrade{5}

% Комментарий по критерию в свободной форме (оставьте пустым, если нет)
\def\continuityComment{}

% Оценка по критерию "Отчётность была регулярной, коммуникации с научным руководителем были не реже раза в неделю", от 1 (совсем нет) до 5 (абсолютно да)
\def\communicationsGrade{5}

% Комментарий по критерию в свободной форме (оставьте пустым, если нет)
\def\communicationsComment{}

% Оценка по критерию "Уровень общей технической грамотности и самостоятельности обучающегося соответствовал ожиданиям, необходимости в разъяснении очевидных вещей не возникало", от 1 (совсем нет) до 5 (абсолютно да)
\def\competenceGrade{5}

% Комментарий по критерию в свободной форме (оставьте пустым, если нет)
\def\competenceComment{}

% Оценка по критерию "Получившийся код соответствует ожиданиям по качеству", от 1 (совсем нет) до 5 (абсолютно да)
\def\codeGrade{5}

% Комментарий по критерию в свободной форме (оставьте пустым, если нет)
\def\codeComment{}

% Оценка по критерию "Замечания своевременно и полностью исправлялись", от 1 (совсем нет) до 5 (абсолютно да)
\def\fixesGrade{5}

% Комментарий по критерию в свободной форме (оставьте пустым, если нет)
\def\fixesComment{}

% Результат работы: 
% pullrequest — должен был быть сделан пуллреквест в существующий репозиторий
% newRepo — должен был быть создан новый репозиторий
% nda — код есть, но закрыт
% что угодно другое, включая пустое значение — кода не предполагалось
\def\contributionType{pullrequest}

% Есть ли внятное описание проекта в README или пуллреквеста в пуллреквесте. "yes", если да, что угодно ещё, если нет.
\def\projectDescription{yes}

% Есть ли модульные тесты. "yes", если да, что угодно ещё, если нет.
\def\tests{yes}

% Есть ли CI с настроенным линтером, в котором проходят и модульные тесты (игнорируется, если речь про пуллреквест). "yes", если да, что угодно ещё, если нет.
\def\ciWithLinter{yes}

% Замерджен ли пуллреквест в основную ветку (игнорируется, если речь про отдельный репозиторий). "yes", если да, что угодно ещё, если нет.
\def\merged{yes}

% Развёрнут ли проект (игнорируется, если речь про пуллреквест). "yes", если да, что угодно ещё, если нет.
\def\deployed{yes}

% Итоговая оценка
\def\grade{A}

% Комментарий к итоговой оценке, если надо. Если нет, оставьте пустым.
\def\gradeComment{}

% Информация о научном руководителе
\def\advisorJobTitle{должность научного руководителя}
\def\advisorScientificDegree{учёная степень, учёное звание}
\def\advisorName{ФИО научного руководителя}

% Генерируемая часть отзыва, тут трогать по идее ничего не надо

\documentclass[a4paper]{article}

\usepackage[a4paper, top=8mm, bottom=8mm, left=8mm, right=8mm]{geometry}

\usepackage{polyglossia}
\setdefaultlanguage[babelshorthands=true]{russian}
\setotherlanguage{english}

\usepackage{fontspec}
\setmainfont{FreeSerif}
\newfontfamily{\russianfonttt}[Scale=1]{DejaVuSansMono}

\usepackage[tiny, compact]{titlesec}

\usepackage{titling}
\setlength{\droptitle}{-1cm}
\pretitle{\begin{center}\begin{bfseries}\Large}
\posttitle{\par\end{bfseries}\end{center}}
\preauthor{\begin{center}\normalsize}
\postauthor{\par\end{center}\vspace{-1.8cm}}

\usepackage{csquotes}
\usepackage{etoolbox}
\usepackage{tabularx}

\newcommand{\checkmarkEntry}[1]{%
    \ifdefstring{#1}{yes}{\def\mark{☑}}{\def\mark{☐}}%
    \mark
}

\title{Отзыв научного руководителя \practiceType\ практике обучающегося\newline
\studentName\ на тему \enquote{\practiceTitle} за \practiceSemester\ семестр}

\author{}

\date{}

\begin{document}

\maketitle

Оценка по шкале от 1 до 5 критериев выполнения практики:

\begin{enumerate}
    \item Цель практики достигнута полностью: \newline
        \textbf{\completenessGrade}\ifdefempty{\completenessComment}{}{, комментарий: \textit{\completenessComment}}.
    \item Этапы работы выполнялись в ожидаемые сроки, работа велась непрерывно в течение семестра: \newline
        \textbf{\continuityGrade}\ifdefempty{\continuityComment}{}{, комментарий: \textit{\continuityComment}}.
    \item Отчётность была регулярной, коммуникации с научным руководителем были не реже раза в неделю: \newline
        \textbf{\communicationsGrade}\ifdefempty{\communicationsComment}{}{, комментарий: \textit{\communicationsComment}}.
    \item Уровень общей технической грамотности и самостоятельности обучающегося соответствовал ожиданиям, необходимости в разъяснении очевидных вещей не возникало: \newline
        \textbf{\competenceGrade}\ifdefempty{\competenceComment}{}{, комментарий: \textit{\competenceComment}}.
    \item Получившийся код соответствует ожиданиям по качеству: \newline
        \textbf{\codeGrade}\ifdefempty{\codeComment}{}{, комментарий: \textit{\codeComment}}.
    \item Замечания своевременно и полностью исправлялись: \newline
        \textbf{\fixesGrade}\ifdefempty{\fixesComment}{}{, комментарий: \textit{\fixesComment}}.
\end{enumerate}

\ifdefstring{\contributionType}{pullrequest}{
    Работа предполагала пуллреквест в существующий репозиторий:
    \begin{itemize}
        \item[\checkmarkEntry{\projectDescription}] есть внятное описание пуллреквеста с технической документацией
        \item[\checkmarkEntry{\tests}] есть модульные тесты на новую функциональность
        \item[\checkmarkEntry{\merged}] изменения приняты в основную ветку
    \end{itemize}
}{
    \ifdefstring{\contributionType}{newRepo}{
        Работа предполагала создание нового репозитория:
        \begin{itemize}
            \item[\checkmarkEntry{\projectDescription}] есть внятное описание проекта в README с техническим описанием, инструкцией по сборке и запуску, примеры использования
            \item[\checkmarkEntry{\tests}] есть модульные тесты
            \item[\checkmarkEntry{\ciWithLinter}] настроен CI с модульными тестами и линтером
            \item[\checkmarkEntry{\deployed}] проект развёрнут и используется
        \end{itemize}
    }{
        \ifdefstring{\contributionType}{nda}{
            Код находится под договором о неразглашении.
        }{
            Программная реализация в рамках практики не предполагалась.
        }
    }
}

Итого рекомендуемая оценка: \textbf{ECTS \grade}.

\ifdefempty{\gradeComment}{}{Комментарий к оценке: \gradeComment.}

\begin{flushright}
    Дата: \today
\end{flushright}

\begin{tabularx}{\textwidth}{@{}X r@{}}
    Научный руководитель: & \underline{\hspace{2cm}}\ \ /\advisorJobTitle,\\
                          & \advisorScientificDegree\ \advisorName. \\
\end{tabularx}

\end{document}